% !TEX root = ../notes.tex

\documentclass[../notes.tex]{subfiles}

\begin{document}

\section{March 2}
Here we go.

\subsection{Modules over Filtered Spectra}
Let's talk some more about filtered spectra.
\begin{definition}[bigraded sphere]
	Let $S^{p,q}$ be the bigraded sphere, which is a filtered spectrum of the form
	\[\cdots\to0\to0\to S^p\to S^p\to\cdots,\]
	where the first $S^p$ appears in degree $2q-p$. (Recall that degrees decrease going down the sequence.) In particular, $q\in\frac12\ZZ$ is permitted.
\end{definition}
\begin{remark}
	All these bigraded spheres are $\otimes$-invertible, where $S^{0,0}$ is the unit.
\end{remark}
\begin{remark}
	With this grading convention, we see that $\Sigma S^{p,q}=S^{p+1,q+1/2}$.
\end{remark}
\begin{example}
	There is a double shift $\tau\colon S^{0,-1}\to S^{0,0}$, which looks like the following.
	% https://q.uiver.app/#q=WzAsMTQsWzQsMCwiU14wIl0sWzUsMCwiU14wIl0sWzYsMCwiXFxjZG90cyJdLFszLDAsIjAiXSxbNCwxLCJTXjAiXSxbNSwxLCJTXjAiXSxbNiwxLCJcXGNkb3RzIl0sWzMsMSwiU14wIl0sWzIsMSwiU14wIl0sWzIsMCwiMCJdLFsxLDAsIjAiXSxbMSwxLCIwIl0sWzAsMCwiXFxjZG90cyJdLFswLDEsIlxcY2RvdHMiXSxbMTIsMTBdLFsxMCw5XSxbOSwzXSxbMywwXSxbMCwxXSxbMSwyXSxbMTMsMTFdLFsxMSw4XSxbOCw3XSxbNyw0XSxbNCw1XSxbNSw2XSxbMSw1LCIiLDEseyJsZXZlbCI6Miwic3R5bGUiOnsiaGVhZCI6eyJuYW1lIjoibm9uZSJ9fX1dLFswLDQsIiIsMSx7ImxldmVsIjoyLCJzdHlsZSI6eyJoZWFkIjp7Im5hbWUiOiJub25lIn19fV0sWzMsN10sWzksOF0sWzEwLDExLCIiLDEseyJsZXZlbCI6Miwic3R5bGUiOnsiaGVhZCI6eyJuYW1lIjoibm9uZSJ9fX1dXQ==&macro_url=https%3A%2F%2Fraw.githubusercontent.com%2FdFoiler%2Fnotes%2Fmaster%2Fnir.tex
	\[\begin{tikzcd}[cramped]
		\cdots & 0 & 0 & 0 & {S^0} & {S^0} & \cdots \\
		\cdots & 0 & {S^0} & {S^0} & {S^0} & {S^0} & \cdots
		\arrow[from=1-1, to=1-2]
		\arrow[from=1-2, to=1-3]
		\arrow[equals, from=1-2, to=2-2]
		\arrow[from=1-3, to=1-4]
		\arrow[from=1-3, to=2-3]
		\arrow[from=1-4, to=1-5]
		\arrow[from=1-4, to=2-4]
		\arrow[from=1-5, to=1-6]
		\arrow[equals, from=1-5, to=2-5]
		\arrow[from=1-6, to=1-7]
		\arrow[equals, from=1-6, to=2-6]
		\arrow[from=2-1, to=2-2]
		\arrow[from=2-2, to=2-3]
		\arrow[from=2-3, to=2-4]
		\arrow[from=2-4, to=2-5]
		\arrow[from=2-5, to=2-6]
		\arrow[from=2-6, to=2-7]
	\end{tikzcd}\]
	Notably, the first $S^0$ in the top occurs in degree $-2$.
\end{example}
\begin{remark}
	If $X_\bullet$ is a filtered spectrum, then the bigraded homotopy groups $\pi_{*,*}X$ is a $\ZZ[\tau]$-module, and with our current grading conventions, one finds that $\tau$ shifts everything down by $2$.
\end{remark}
Recall that we are interested in the filtered $\mathbb E_\infty$-ring
\[\op{fil}^*\mathbb S_{(p)}=\lim_\Delta\tau_{2\ge*}\op{BP}^{\otimes(\bullet+1)}.\]
Taking the cofiber by $\tau$ passes to the associated graded, so one finds that
\[\pi_{2t-s,s}(\op{fil}^*\mathbb S/\tau)=\mathrm H^s\left(\mc M_{\mathrm{fg},(p)}^{\mathrm{strict}}\right),\]
which is also $\mathrm H^s(\mc M_{\mathrm{fg},(p)};\omega^{\otimes t})$. (Recall that the formal groups up to strict isomorphism have an action by $\mathbb G_m$, which is where the extra grading comes from on the right-hand side.) On the other hand, inverting $\tau$ amounts to taking the colimit of the filtration, so we see that
\[\pi_{**}\op{fil}^*\mathbb S_{(p)}\left[\tau^{-1}\right]=\pi_*\mathbb S_{(p)}.\]
It now turns out that $\tau^r$-torsion elements correspond to $d_{2r+1}$-differentials: in each graded piece, there will be one element in the kernel and one element in the cokernel, which are the source and target for our differential.\footnote{The sort of example to have in mind is that $\tau\colon\ZZ[\tau]/\tau^r\to\ZZ[\tau]/\tau^r$ has one element in its kernel and cokernel.} In particular, if an element is going to live all the way to $\pi_*\mathbb S_{(p)}$, then it is not going to be $\tau$-torsion and instead goes all the way down the complex.
\begin{remark}
	Note that $\op{fil}^*\mathbb S/\tau$ is an $\mathbb E_\infty$-ring: moving around the quotient, it is $\op{fil}^*\mathbb S\otimes\left[S^{0,0}/\tau\right]$, and the latter is an $\mathbb E_\infty$-ring as the tensor product of $\mathbb E_\infty$-rings. We can thus consider modules over $\op{fil}^*\mathbb S/\tau$, which turns out to be the derived category over $\mc M_{\mathrm{fg},(p)}$. (The idea is that modding out by $\tau$ passes to the associated graded, which remains true ``on the level of categories.'') Notably, the elements of this category are filtered colimits of bounded chain complexes of $\op{BP}_*\op{BP}$ comodules. In particular, we have a large source of objects given by $\op{BP}_*\op{BP}$-comodules.
\end{remark}
\begin{example}
	One has a comodule $\mathrm{BP}_*$ over the Hopf algebroid $\op{BP}_*\op{BP}$. Similarly, we have $\op{BP}_*/p$.
\end{example}
\begin{nex}
	Recall that we can present $\op{BP}_*$ as $\ZZ_{(p)}[v_1,v_2,\ldots]$. Then $\op{BP}_*/v_1$ fails to have a comodule structure because the right map $\mathrm{BP}_*\to\mathrm{BP}_*\mathrm{BP}$ is given by $v_1\mapsto v_1+pt_1$, so we cannot ``just'' take a quotient by $v_1$. However, the quotient $\op{BP}_*/(p,v_1)$ is okay.
\end{nex}
\begin{example}
	In general, one has that $\eta_R(v_i)\equiv v_i\pmod{p,v_1,\ldots,v_{i-1}}$, so we are able to form the quotients $\mathrm{BP}_*/(p,v_1,v_2,\ldots)$. Notably, this is just $\FF_p$.
\end{example}
\begin{example}
	If $X$ is any spectrum, then we have a comodule $\op{BP}_{2*}X$. On homotopy, this looks like
	\[\pi_{2*}(\mathrm{BP}\otimes X)\to\pi_{2*}(\mathrm{BP}\otimes\mathrm{BP}\otimes X)\to\cdots,\]
	so the action by $\op{BP}_*\op{BP}$ is a little clearer.
\end{example}
In general, if we are handed a module over $\op{fil}^*\mathbb S/\tau$, we can ask it comes from a module $M$ over $\op{fil}^*\mathbb S$. In particular, $M\left[\tau^{-1}\right]$ is then a spectrum which ``lifts'' its quotient $M/\tau$.
\begin{remark}
	Here is an open question: fix $n$, and take $p$ large. Then we do not know if the module
	\[\frac{\op{BP}_*}{(p,v_1,\ldots,v_n)}\]
	lifts to a spectrum. For example, even at $n=1$, this is not possible if $p=2$. This was recently proved for $n=4$.
\end{remark}
To continue, we need more module theory.
\begin{definition}[thick subcategory]
	Fix an $\mathbb E_\infty$-ring $R$. Given a collection $\{M_i\}_{i\in I}$ of $R$-modules, the \textit{thick subcategory} of $R$-modules generated by these modules is the smallest full subcategory containing the modules and closed under finite limits, finite colimits, and retracts.
\end{definition}
\begin{definition}[perfect]
	Fix an $\mathbb E_\infty$-ring $R$. Then a \textit{perfect} module is an $R$-module in the thick subcategory generated by $R$.
\end{definition}

\subsection{Krause's Theorem}
By extending the methods from last class, we can show the following.
\begin{theorem}[Krause] \label{thm:krause}
	Fix a perfect module $X$ over $\op{fil}^*\mathbb S_{(p)}$. Then there is a slope
	\[s\in\left\{\frac1{p^{j+1}\left(p^i-1\right)}:i>j\ge0\right\}\]
	satisfying the following properties.
	\begin{listalph}
		\item $\pi_{*,*}X$ s concentrated below a line of slope $s$, and this slope $s$ is minimal with this property.
		\item There is an endomorphism $\Theta$ of $X$ with that slope such that $\op{cofiber}\theta$ admits a smaller slope.
	\end{listalph}
\end{theorem}
\begin{example}
	Recall that there is a functor $\nu$ from spectra to filtered spectra which sends $X$ to
	\[\nu(X)\coloneqq\lim\tau_{2*}\left(X\otimes\mathrm{BP}^{\otimes(\bullet+1)}\right).\]
	If $X$ is perfect, then $\nu(X)$ is perfect.
\end{example}
\begin{remark}
	For any perfect $X$ over $\op{fil}^*\mathbb S_{(p)}$, the slope of $X$ agrees with the slope of $X/\tau$. Indeed, $\tau$ simply moves all the homotopy groups down, so slopes from $X$ descend to $X/\tau$. It also turns out that slopes for $X/\tau$ come from algebra (which follows from the proof) and thus lift to $X$.
\end{remark}
\begin{proof}[Idea]
	By modding out by $\tau$, we can reduce this to an algebra statement. In particular, for a perfect module $M$ over $\op{fil}^*\mathbb S/\tau$, we can form the quotient $M/(p,v_1,v_2,\ldots)$, and it is enough to prove the theorem for the slopes of this quotient.

	However, the $v_i$s basically shift the elements in $\pi_{**}M$ to the right. Let's now explain how this relates to what we discussed last class. Note that we can form the quotient $\mathrm{BP}_*\mathrm{BP}/(p,v_1,v_2,\ldots)=\FF_p[t_1,t_2,\ldots]$. On the other hand, note that
	\[\pi_*\left(\frac{\mathrm{BP}\otimes\mathrm{BP}}{(p,v_1,v_2,\ldots)}\right)=\pi_*(\FF_p\otimes\mathrm{BP})\to\pi_*(\FF_p\otimes\FF_p)^m,\]
	where the second map is simply given by projecting $\op{BP}_*$ to $\op{BP}_0$ via the Whitehead tower, which is $\FF_p$. Now, this right-hand algebra is $\FF[t_1,t_2,\ldots]\otimes\land^\bullet(\sigma t_1,\sigma t_2,\ldots)$, and it is the sort of thing we have some control over from last class: it is an iterated extension of some $\FF_p[x]/\left(x^{p^r}\right)$s.
\end{proof}
\begin{remark}
	One can view \Cref{thm:krause} as a ``purely algebraic'' result. Indeed, it boils down to understanding the diagram
	% https://q.uiver.app/#q=WzAsMyxbMCwwLCJcXG1hdGhybXtNb2R9KFxcb3B7ZmlsfV4qXFxtYXRoYmIgUykiXSxbMSwwLCJcXG1jIEQoXFxtYyBNX3tcXG1hdGhybXtmZ30sKHApfSkiXSxbMCwxLCJcXG1hdGhybXtTcGVjdHJhfSJdLFswLDIsIlxcdGF1XnstMX0iLDJdLFswLDEsIi9cXHRhdSJdXQ==&macro_url=https%3A%2F%2Fraw.githubusercontent.com%2FdFoiler%2Fnotes%2Fmaster%2Fnir.tex
	\[\begin{tikzcd}[cramped]
		{\mathrm{Mod}(\op{fil}^*\mathbb S_{(p)})} & {\mc D(\mc M_{\mathrm{fg},(p)})} \\
		{\mathrm{Spectra}_{(p)}}
		\arrow["{/\tau}", from=1-1, to=1-2]
		\arrow["{\tau^{-1}}"', from=1-1, to=2-1]
	\end{tikzcd}\]
	as well as the full understanding of the category $\mc D(\mc M_{\mathrm{fg},(p)})$ (along with their slopes) produced last class.
\end{remark}

\subsection{The Main Theorem}
We are now ready to prove \Cref{thm:slope}.
\slopethm*
\begin{proof}
	By \Cref{thm:krause}, there is a sequence of $\op{fil}^*\mathbb S$-modules given by successively modding out by the produced endomorphisms $\Theta_1,\Theta_2,\ldots$. By inspection, any decreasing sequence in our set of slopes must go to zero. We can refine this choice of endomorphisms using Nishida nilpotence.
	\begin{example} \label{ex:theta-1-tau-torsion}
		Consider $\pi_{**}\op{fil}^*\mathbb S_{(2)}$. Then our homotopy groups feature $\eta$ in degree $1$ which is not $\tau$-torsion, and $\eta$ would suffice to provide the endomorphism in slope $1$. Additionally, $\eta^2$ is not $\tau$-torsion, nor is $\eta^3$, but $\eta^4$ does vanish under $\tau$. Thus, it is possible to choose this endomorphism $\Theta_1$ to be $\tau$-torsion.
	\end{example}
	\begin{lemma}[Birkland]
		Let the decreasing slopes of $\op{fil}^*\mathbb S$ from \Cref{thm:krause} be $\{s_1,s_2,\ldots\}$. Then one can choose associated endomorphisms $\Theta_1,\Theta_2,\ldots$ of the correct slopes so that each $\Theta_i$ is $\tau$-power torsion.
	\end{lemma}
	\begin{proof}
		\Cref{ex:theta-1-tau-torsion} explains how to choose $\Theta_1$. Let's explain how we could choose $\Theta_2$, and the remaining $\Theta$s are basically chosen by the same process. Then we are on the hunt from some endomorphism of $\op{fil}^*\mathbb S/\Theta_1$ (up to the usual shifts giving the slope $s_2$). After inverting $\tau$, we see that $\Theta_1$ vanishes because it is $\tau$-torsion, so taking this cofiber produces $\mathbb S\oplus\mathbb S$ (again, up to some suspensions).

		Now, suppose that we use \Cref{thm:krause} to produce some endomorphism $\Theta_2$. Then the key trick is that \Cref{thm:nishida-nilpotence} tells us that $\tau^{-1}\Theta_2$ is a $2\times2$ matrix whose entries are nilpotent. It follows that $\tau^{-1}\Theta_2$ is nilpotent, so $\tau^a\Theta_2^b=0$ for some $a$ and $b$. Thus, we simply replace $\Theta_2$ with $\Theta_2^b$, and we are done!
	\end{proof}
	\begin{lemma}
		Choose $\tau$-power torsion endomorphisms $\Theta_1,\Theta_2,\ldots$ of $\op{fil}^*\mathbb S$. For each nonnegative integer $n$, the module $\op{fil}^*\mathbb S$ is in the thick subcategory of $\op{fil}^*\mathbb S$-modules generated by
		\[\left\{\frac{\op{fil}^*\mathbb S}\tau,\frac{\op{fil}^*\mathbb S}{\Theta_1,\ldots,\Theta_n}\right\}.\]
	\end{lemma}
	\begin{proof}
		We proceed by induction; for $n=0$, there is nothing to do. For $n\ge1$, consider the cofiber sequence
		\[\frac{\op{fil}^*\mathbb S}{\Theta_1,\ldots,\Theta_n}\to\frac{\op{fil}^*\mathbb S}{\Theta_1,\ldots,\tau\Theta_n}\to\frac{\op{fil}^*\mathbb S}{\Theta_1,\ldots,\Theta_{n-1},\tau}.\]
		It follows that we can build the middle term from the two outer terms, so the middle term lives in the thick subcategory. Iterating this process allows us to build the module
		\[\frac{\op{fil}^*\mathbb S}{\Theta_1,\ldots,\tau^a\Theta_n}\]
		for any $a$, so sending $a$ large completes the proof by our induction.
	\end{proof}
	Now, anything in the thick subcategory generated by these two modules admits a line of slope coming from the slope of $\op{fil}^*\mathbb S/(\Theta_1,\ldots,\Theta_n)$, where everything above the line is $\tau$-power torsion. It follows that the actual homotopy groups (on $E_\infty$) must always live below this line! Letting $n$ go to infinity, we see that the slopes go to zero. This completes the proof of \Cref{thm:slope}.
\end{proof}

\end{document}